\documentclass[a4paper, 12pt]{book}
\usepackage{amssymb}
\usepackage{xspace}
\usepackage{tikz}
\usepackage[newitem,newenum,neverdecrease]{paralist}
\usepackage{amsmath,amssymb,amsfonts,mathrsfs,latexsym,stmaryrd}
\usepackage{mathtools}
\usepackage{amsthm}
\usepackage{ifpdf}
\usepackage{cases}
\usepackage{ragged2e} 	% https://tex.stackexchange.com/questions/89680/how-can-one-set-full-justification-within-left-justified-raggedright-text
\usepackage{breakcites} % https://tex.stackexchange.com/questions/2773/how-do-i-make-latex-push-long-citations-to-a-new-line
\usepackage{hyperref}	% http://tex.stackexchange.com/questions/73862/how-can-i-make-a-clickable-table-of-contents
\usepackage{fancyhdr}	% https://www.sharelatex.com/blog/2013/08/06/thesis-series-pt2.html
\usepackage{titlesec} 	% http://tex.stackexchange.com/questions/11444/how-to-format-the-chapter-heading
\usepackage[T1]{fontenc}
\usepackage[utf8]{inputenc}
%\usepackage{setspace}
\usepackage{xcolor}
\usepackage{geometry}
\usepackage{layout}		% http://tex.stackexchange.com/questions/50258/margins-of-book-class
\usepackage{listings}
\usepackage{booktabs}
\usepackage{graphicx}
\graphicspath{ {./img/} }

\usepackage{pseudocode}
\usepackage{chngcntr}


% fixes for pseudocode-hyperref interaction
\counterwithout{pseudocode}{section}
%\counterwithin{pseudocode}{chapter}
\renewcommand{\thepseudocode}{\thechapter.\arabic{pseudocode}}

\newcounter{newpseudonum}[pseudocode]
\renewcommand{\thenewpseudonum}{\thepseudocode\alph{newpseudonum}}
\providecommand{\labelline}[2][0.5in]{&\refstepcounter{newpseudonum}\hspace*{#1}\mbox{(\thenewpseudonum)}\label{#2}}
\ifpdf
  \providecommand{\refline}[1]{\hyperref[#1]{(\ref*{#1})}}
\else
  \providecommand{\refline}[1]{\ref*{#1}}
\fi
\newcommand{\pcodetab}[1]{\hspace*{#1ex}}

% alignement fixes for pseudocode blocks
\renewcommand{\DO}{\\\!\!\pcodetab{1}\mbox{ \bfseries \makebox[0pt][l]{do}\phantom{xx} }}
\renewcommand{\ELSE}{\\\pcodetab{1}\mbox{ \bfseries \makebox[0pt][l]{else}\phantom{then} }}
\renewcommand{\RETURN}[1]{\ifthenelse{\equal{#1}{} }{\mbox{\bfseries return}}{\mbox{\bfseries return}#1}}
\newcommand{\FUNCTION}[2]{\mbox{\bfseries proc }\mbox{\textsc{#1}}\left(\ensuremath{#2}\right)\\}
\newcommand{\ENDFUNCTION}{}
\newcommand{\FCALL}[2]{\mbox{\textsc{#1}}\left(\ensuremath{#2}\right)\\}



% code box length
\newlength{\pcodewidth}
\setlength{\pcodewidth}{\textwidth}
\addtolength{\pcodewidth}{-58pt}

% code environment
\newenvironment{code}[1]{
\begin{Sbox}
\!\!\begin{minipage}{#1}%\pcodewidth}
\bfseries
\noindent
%\begin{math}
\small
$$
\begin{array}{@{\hspace*{1ex}}lr@{}}
}{
\end{array}
$$
%\end{math}
\end{minipage}\vspace{-2mm}
\end{Sbox}
\shadowbox{\TheSbox}{}
}	% for pseudocodes

%%% Formattazione del header del Capitolo
\titleformat
{\chapter} % command
[display] % shape
{\bfseries\Huge\itshape} % format
{\flushright \color{black!45}{\Large Chapter \thechapter}} % label
{0.5ex} % sep
{
    \rule{\textwidth}{3pt}
    \vspace{1ex}
    \centering
} % before-code
[
\vspace{-0.5ex}%
\rule{\textwidth}{3pt}
] % after-code
%%% end

% https://tex.stackexchange.com/questions/95488/list-of-figures-and-page-numbering
\makeatletter
\newcommand{\emptypage}[1]{%
  \cleardoublepage
  \begingroup
  \let\ps@plain\ps@empty
  \pagestyle{empty}
  #1
  \cleardoublepage}
\makeatletter

% base line strech (default 1.0) -- interlinea
\renewcommand{\baselinestretch}{1.2}

\begin{document}
%%% Nuova geometria per la pagina del titolo
% visto che la classe del documento è book, le pagine pari e dispari avranno geometrie diverse,
% mentre quella del titolo deve essere unica!
\newgeometry{
  top=2cm,
  bottom=2.5cm,
  left=2.5cm,
  right=2.5cm,
  headsep=25pt,
  headheight=14.5pt
}

%%% Pagina del titolo
\begin{titlepage}
%	\topskip0pt
	%\vspace*{\fill}
	\centering
%	\vspace*{20mm}
	\includegraphics[]{logo.png}\\
	\vspace*{1cm}
	\huge \textbf{\textsc{Università degli Studi di Ferrara}}
	\Large \textsc{Corso di Laurea in Informatica}
	
	\vspace*{1.5cm}
	\hrule width \hsize \kern 1mm \hrule width \hsize height 2pt
	\vspace*{10mm}
	\Huge \emph{\textbf{Data diodes for secure Industry 4.0 networking: \\ A railway environment monitoring use case}}
	\vspace*{10mm}
	\hrule width \hsize height 2pt
	\vspace*{1mm}
	\hrule width \hsize \kern 1mm
	
	\vspace*{5mm}
	\begin{minipage}{0.46\textwidth}
		\begin{flushleft} \Large
			\emph{Relatore:}\\
			\Large \textbf{Prof. Carlo \textsc{Giannelli}}
            \emph{Tutor Aziendale:}\\
			\Large \textbf{Ing. Roberto \textsc{Giansante}}
		\end{flushleft}
	\end{minipage}
	\begin{minipage}{0.45\textwidth}
		\begin{flushright} \Large
			\emph{Laureando:} \\
			\Large \textbf{Giuseppe \textsc{Cervone}}
		\end{flushright}
	\end{minipage}
	
	\vspace*{20mm}
	\Large \textsc{Anno Accademico $2020-2021$}
\end{titlepage}
\restoregeometry

\addtocontents{toc}{~\hfill\textbf{Page}\par}	% https://texblog.org/2011/09/09/10-ways-to-customize-tocloflot/
\pagestyle{empty}
\clearpage
\tableofcontents
\thispagestyle{empty}
\addtocontents{toc}{\protect\thispagestyle{empty}}	% http://tex.stackexchange.com/questions/2995/removing-page-number-from-toc


\chapter{Riassunto}

In questo lavoro di tesi si andranno a discutere i concetti di Industria 4.0, Industrial
Internet of Things e Cybersecurity. Verranno utilizzate le definizioni di questi concetti
per descrivere il flusso lavorativo della sede di Bologna di Alstom, compagnia per cui
è stato svolto il lavoro di tirocinio esposto in questa tesi. Alstom è un’azienda
multinazionale che si occupa di produzione e di monitoraggio, sia delle vie
ferroviarie, che dei veicoli che viaggiano su di esse. In particolare, la sede di
Bologna viene considerata un centro di eccellenza per quello che è il monitoraggio.
Si andrà quindi a discutere del pensiero e della struttura della soluzione studiata al
fine di risolvere il problema posto da Alstom, dei singoli componenti che ne fanno
parte e di come è possibile ricreare la suddetta soluzione in caso si volesse testarne
le funzionalità e, infine, si descriverà anche degli eventuali miglioramenti identificati
che potrebbero essere applicati in futuro.
La soluzione studiata, in dettaglio, si pone come obiettivo quello di tentare di
replicare il funzionamento di un diodo dati a basso costo. Un diodo dati è un
dispositivo di uso comune nel settore della cyber security, che permette la
comunicazione unidirezionale di dati da una rete protetta verso una rete pubblica.
Questo dispositivo offre quindi l’opportunità a chi lo utilizza, di ottenere informazioni
sui dati presenti in una rete privata, senza la necessità di instaurare un canale di
comunicazione diretto con essa. Il diodo dati in questo caso verrà simulato grazie
alla comunicazione seriale tramite lo standard RS-232, comunicazione che viene
resa unidirezionale tagliando il terminale responsabile per l’invio di dati verso il lato
privato della comunicazione. Il cavo seriale in questione verrà connesso a due
microcomputer, in dettaglio due Raspberry Pi 3. Questi ultimi saranno responsabili
della sincronizzazione nella comunicazione e di rendere le informazioni comunicate
facilmente da reperibili. Questo processo verrà svolto tramite l’ausilio di Zabbix, un
software di monitoraggio per reti che offre anche la possibilità all’utente di simulare
un ambiente con dati provenienti da altri fonti. Saranno quindi poste in esecuzione
due istanze di Zabbix, la prima che sarà dislocata sulla Raspberry identificata come
“privata”, e che catturerà i dati dai dispositivi nella rete privata, e l’altra che andrà ad
essere collocata sulla Raspberry definita "pubblica", la quale riceve i dati tramite
comunicazione unidirezionale. Questo scenario andrà a creare sulla seconda
instanza di Zabbix, menzionata in precedenza, una replica di quello che succede
realmente sulla rete privata e protetta.


\chapter{Introduction}
%%% Fancy header settings, queste impostazioni vanno fatte solo una volta all'inizio del primo capitolo!
\pagestyle{fancy}
\fancyhf{}
\renewcommand{\headrulewidth}{2pt}
\fancyhead[EL]{\textbf{\textsf{\nouppercase\thepage}}}
\fancyhead[ER]{\textbf{\textsf{\nouppercase\leftmark}}}
\fancyhead[OR]{\textbf{\textsf{\nouppercase\thepage}}}
\fancyhead[OL]{\textbf{\textsf{\nouppercase {\rightmark}}}}
%%% end

Industry 4.0, cybersecurity, Internet Of Things, these are all terms that 30 years ago,
when the idea of “Internet” first came into our lives, would have been considered
gibberish, but that now commonly used in today’s society, but how did they come to
be? What are they exactly? And how can they improve our lives? These and many
other questions are what I will try to tackle with my thesis work, looking at their
advent, the modern perception of them, how these concepts are changing the
industrial space, and how it is possible to build something by putting all these
concepts together to make someone’s life a little easier.
I will do so by dissecting in the first chapter of this thesis the concept of Industry 4.0
and how it’s advent led to Internet of Things devices becoming ever more prevalent
in our day to day use. This will also help us understand why the idea of cybersecurity
is not to be ignored. Then in chapter two, I will look at the thought process of a
multinational company, trying to understand their workflow, and how the researching
of new solutions to improve their services comes to fruition. These new solutions
obviously are compliant with the concepts mentioned earlier. With the information
gathered about the company, and with full knowledge of the new possibilities given
to us by the enabling technologies of Industrial IoT, I will dissert in the third chapter
how a IIoT solution for a multinational is built. In my particular case it will be a
solution with a microprocessor, with interconnectivity between devices, and with the
possibility of monitoring devices under a network.

%%% all'inizio di ogni capitolo, questa impostazione rimuove il numero di pagina, provare a commentare per vedere la differenza
\thispagestyle{empty}

\newpage

\chapter{Industry 4.0 and cybersecurity}
A common practice in studying history is ‘Periodization’. The concept behind it is that dividing historical periods in well defined time frames, allows us to refer more easily to a period of time and allows us to better understand the connections that current events might have to events which happened in the past [IAT14]. An example of such periodization, which uses a timeframe, events linked to that specific timeframe and the idea that time periods can be linked using the events happening in that time frame, is the periodization of Industrial Revolutions.
The Industrial Revolutions, which have been divided in four different timeframes, are a  perfect example of the practice of Periodization put to use. Having four different time periods being referred to as a period of revolution in industrial processes, highlights connections in the events happening between each one of those time frames in history.
The connection between each of the Industrial Revolutions is that in each one of these time periods, introductions of new revolutionary technologies occured in the manufacturing space. The First Industrial Revolution, which started in the 18th century, is characterized by the introduction of the first machinery which substituted hand production methods for some of the tasks. The Second Industrial Revolution, usually linked to World War I, saw the widespread adoption of technological systems such as telegraphs and railroad networks amongst other introductions such as assembly lines. The Third Industrial Revolution, which began in the 70’s saw the introduction of the first automated tasks and the first computers introduced in the assembly line. The Fourth Industrial Revolution begins in the 2010s, is the revolution we are currently living, and the one we will be focusing on [IRF].

\section{The Industry 4.0 movement}
The Fourth Industrial Revolution, now commonly known as Industry 4.0 is characterized as the time period in history where implementation of information and communication technologies is taking over traditional manufacturing practices. While the last two Industrial Revolutions both heavily involve the introduction of computers in the manufacturing space, they are differentiated by the fact that in the Fourth Industrial Revolution, the interaction between humans and technology is much more closely knit, leading to one unified network, and thus to the creation of “cyber-physical production systems” [IRF].
As mentioned earlier, the Fourth Industrial Revolution is often referred to as the Industry 4.0 movement. The terminology comes from Germany, first seen in 2011, the term ‘Industrie 4.0’ was used by Henning Kagermann, Wolf-Dieter Lukas and Wolfgang Wahlster, in a paper called “Industrie 4.0: Mit dem Internet der Dinge auf dem Weg zur 4 industriellen Revolution”, which translates to “Industry 4.0: Internet of Things on the Road to the Fourth Industrial Revolution”. In this paper, which was presented at the Hannover Fair, they presented a set of investments to the German government going from infrastructure and even touching public schools, investments which were meant to propel German manufacturing back to a state of dominance in the global economy [I4M11].
Germany definitely wasn’t the only country that realized a shift in manufacturing practices was needed to keep up with the demands of the market. For example, also in North America another project was initialized to modernize manufacturing technologies. Not necessarily called Industry 4.0, the term used was Industrial Internet, and it was used by the General Electric company in late 2012. The term was used in the foundation of the Industrial Internet Consortium. The consortium is open to any company, as it was founded to accelerate the development, adoption and widespread use of interconnected machines, devices and intelligent analytics. The consortium was later rebranded in August 2021 as Industry IoT Consortium, moving it’s main goal to delivering business value to industry, organizations, and society by accelerating adoption of trustworthy IoT technologies. It is worth mentioning that the aim of the consortium is not to sell products or services but to drive progress and to help members get the best return on their IoT investments [ICB17].
Industry 4.0 thus is a very recent concept, therefore it has different interpretations based on the needs of the manufacturing industry of each nation. There have been attempts however at trying to standardize what being an Industry 4.0 compliant technology means. An example of one of the most cited attempts has been written by Mario Hermann, Tobias Pentek and Boris Otto in a Hawaii International Conference on System Sciences (HICSS) research paper. In this paper, they divided Industry 4.0 principles into four main categories:

\begin{itemize}
    \item interconnection, allowing machinery, sensors and other devices to be connected with each other, and to be connected with the people monitoring them;
    \item information transparency, thus the ability to collect a lot more diagnostic data on each machine, so as to further analyze and optimize current processes;
    \item decentralized decisions which leads to increased automation of tasks. With the inclusion of more diagnostic data being captured, it is easier to optimize tasks even from outside of the production facilities;
    \item technical assistance, with systems assisting humans in decision-making and the ability of systems to perform tasks that are less safe for humans. Since the systems collect more diagnostic data, they could also be more aware of issues, helping users solve them or even automatically solve tasks without need for human intervention [DPI16].
\end{itemize}

These four main design principles most likely derive from looking at the development and improvement of a series of enabling technologies for Industry 4.0. The definition of enabling technology is: ”an invention or innovation that can be applied to drive radical change in the capabilities of a user or culture”. Thus the development that those technologies received in recent times helped the widespread application in the manufacturing space. These enabling technologies have been identified by the Boston Consulting Group in 2015 and have been divided as follows:

\begin{itemize}
    \item advanced manufacturing solutions thanks to robots becoming autonomous, flexible and cooperative;
    \item additive manufacturing, the classic example of which is 3D printing, enabling companies to create customized products for their manufacturing line;
    \item augmented reality systems, which have a variety of uses in the manufacturing space. For example easily being able to place virtual machines in a real environment is a move that if done before building and placing new machinery could dramatically cut costs;
    \item simulations in all stages of production, useful to mirror the real world thus increasing quality, decreasing costs and making production more efficient;
    \item improved horizontal and vertical system integration, meaning that all departments of a business are much more closely linked thanks to data being shared;
    \item industrial internet of things, field devices communicate and interact with more centralized controllers, thus devices, which didn’t have internet, now have it;
    \item the cloud, moving machine data and functionality to the cloud is achieving levels closer to local storage and computability;
    \item cyber security, due to the more common interconnection of devices, protecting critical industrial systems is a big focus of current manufacturing systems;
    \item big data analytics. Analytics based on large data sets optimizes production quality, saves energy and improves equipment service [BCG15].
\end{itemize}

The type of products, services and machines which are compliant with the guidelines described above can be further grouped in these macro-categories: Cyber-physical systems, Internet of Things devices or platforms, Cognitive computing applications and on-demand availability of computer system resources [HDI17].
In this thesis, we will be analyzing further the Internet of Things macro-category. We will be looking at where Industrial Internet of Things as a field comes from, it’s widespread implementation of its technologies in various fields of manufacturing, and some of the limitations which may arise with its inherent flaws.

\section{Industrial Internet of Things}
Industrial Internet of Things is an obvious ramification of the more well known idea of Internet of Things. As with Industry 4.0, origins of the term Internet of Things are many, the most probable of which is credited to Kevin Ashton in 1999. He used the term in a presentation for Procter and Gamble, introducing to them the idea of adding the then new technology of RFID in the supply chain. Ashton’s idea of IoT is although slightly different from what IoT is perceived to be nowadays, Ashton thought that people often went unnoticed in the idea of a network of interconnected devices, but they were actually vital as they are the main provider of information. However, if computers are empowered with ways of gathering information about themselves and their environment, they can capture more information than humans, without human error, and can analyze the captured data. This combination of factors lets us delegate diagnostic work to computers thus allowing humans to concentrate on different parts of the manufacturing process [RFI09].
The evolution of Internet of Things, as it’s more commonly known today, has standardized a new definition of the concept, more inclined towards the idea of embedded objects with sensors, processing ability, software and other technologies that allow for them to connect and interchange data between devices via the internet, also giving internet connectivity to devices that in the past were unable to do so [WII21]. The peculiarity behind the modern concept of IoT, is that while the original concept stemmed from the industrial space, nowadays IoT has been massively applied in the consumer electronics space so much so, that the definition of IoT itself changed. Applications of IoT concepts are now common in household appliances, this can be seen from the fact that fridges, ovens, dishwashers, coffee machines, vacuum cleaners and light bulbs nowadays are all connected to the internet, they are all integrated to voice assistants, they all have their own mobile applications for monitoring and controlling, slowly becoming the norm in modern households.
While IoT is now common in households, we want our focus to stay on the sector that is now segmented by the name Industrial Internet of Things. Historically, control and automation of manufacturing processes was first introduced with the use of Programmable Logic Controllers (PLC) solutions. PLCs, invented by Dick Morey in 1968, are basic microprocessors built with the idea of being long lasting and easily programmable to allow multiple input/output arrangements in real time. The possibility of making the machine work differently based on their input is what makes it the first example of automation in industrial processes [MAM08].
The most obvious evolution of PLCs can be seen in Distributed Control Systems (DCSs). First introduced by Honeywell and Yokogawa in 1975, they improve on PLCs since they allow for automation and control over bigger sets of machinery, distributing control across the whole plant rather than being centralized thus removing the idea of single point of failure. The concept of having a control room with information about the plant, also evolved thanks to the adoption of DCSs, as the technology allows for high levels of plant status and production monitoring. The added set of functionalities brought the drawback that DCSs are slightly harder to program, and slightly less reactive when given an input, since further processing of information is involved. These two aforementioned solutions are not the only examples of industrial automation which came before adoption of IIoT solutions. After the advent of PLCs and DCSs, a variety of similar control systems for factories such as SCADA also gained in popularity over time. SCADA, the acronym of Supervisory Control And Data Acquisition, is a system of software and hardware components built specifically for monitoring and automation of industrial machinery. SCADA systems are often employed in mission critical contexts, meaning areas where risk of failure could lead to extremely dangerous consequences such as gas and oil distribution. All these various control systems differentiate from one to the other although still offering a set of similar features, furthermore they are mostly closed solutions to one single manufacturer. This in turn means that the invention of a more modern automation solution never fully makes the application of an older solution obsolete, even at times integrating elements of the older solution in the newer solution [EIC13].
Focusing back on Industrial IoT, its solutions can be seen as an evolution, and not as another variation on DCS solutions. This differentiation is done because of a series of enabling technologies, the implementation of which allows for clearer nomenclature. Initially the advancements in cloud computing brought massive flexibility to monitoring manufacturing processes remotely, since it is possible to diagnose issues without being on site, making the technicians jobs much easier. The general advancement in hardware, due to the reduced cost and size of microcontrollers and microprocessors, leads generally to the technology being more applicable in manufacturing scenarios compared to proprietary devices. Often, these microcontrollers and microprocessors, can even be off the shelf products, fully programmable to one’s needs and with open source support widely available. Another advancement which brought to standardization of IIoT solutions are the new and well refined standards for wireless communication, both via the web and via solutions such as Bluetooth Low Energy for short-range communications. On the web side, publish/subscribe protocols such as MQTT or request/response RESTful protocols are both valid alternatives based on the use-case scenario to exchange information between machinery.
Having talked about single advancements in IoT and thus IIoT technologies, it is now possible to identify what an IIoT solution is and the various building blocks that create an IIoT solution. The building blocks of an IIoT solution are also defined as the links in a value chain. The value chain illustrates how the different building blocks operate, and what each block contributes to the overall experience of using an IIoT solution. An illustration of the IIoT value chain can be seen in Figure 1.1, where we can see the building blocks listed in the following way:

\begin{itemize}
    \item data capturing module. The data can be either sensor or actuator data, or data caught from a software about machinery;
    \item preprocessing and pre analyzing module. Thanks to edge computing, a modern distributed computing approach in which bandwidth usage is massively cut, the data is initially categorized so as to make later analysis more relevant;
    \item processing data module, where the caught data is actually processed, and where automation commands occur;
    \item data analysis module, where machine learning, artificial intelligence, and big data techniques are used to analyze the data, so as to further improve data processing and data capturing algorithms.
\end{itemize}

As we can also see from Figure 1.1, these modules are also built to interact and communicate the results with each other [IVC20]. The interoperability of modules in IIoT solutions is one of it’s main selling points, however this interoperability is one that also raises the most common concern with IIoT solutions, how secure is a system with openly interchangeable pieces?

\section{Cybersecurity and Industrial IoT}
In cybersecurity there’s a saying: “Fast, secure, cheap, pick two“. This saying is related to the self-explanatory idea that it is often difficult for these three features to coexist in computer science based solutions. This triangle of options is often true because of a series of factors, the most common of which is that the standards on which the internet relies on, were created in a time where cybersecurity wasn’t a vital part of design, but only an afterthought. Since Industrial IoT solutions can also be considered computer science solutions, the concept stands true also in this scenario. This raises the obvious issue of which of the three is better to cut, and the answer is that it depends. For example, in the case of communications where no sensible data is shared, it is often considered ok for the communication to be less secure, on the other hand sensible data can in some cases be shared without hurry, so speed is sacrificed. Convenience in communication is often lost in the case of devices from different manufacturers communicating between different blocks, while it is true that Industrial IoT allows for flexibility between layers as protocols are standard, the added flexibility is not always free, thus sacrificing inexpensiveness.
A different cybersecurity issue altogether, not strictly related to the Industrial IoT space but still worth mentioning, is the issue of well known manufacturers that have never built computer science based solutions. It is an issue that is much more common in the IoT space for households in house appliances, where said manufacturers are usually not educated on the concept of cybersecurity itself, often completely omitting its principles when creating hardware or software for IoT.
While the security concerns might deter companies from investing in Industrial IoT solutions, there are hardware solutions easily applicable to IIoT solutions that are often used, and that haven’t been mentioned yet. The one we will be talking about is a specific hardware solution that can help with making data transfer secure, it does so by making data transfer physically unilateral, securing certain pieces in a closed network, with those pieces outputting data to the end user. This hardware is commonly called a data diode, and as mentioned earlier, it is used as an unilateral single point of communication to send data from an isolated network towards a public network.
Data diodes, as we can see from the diagram in Figure 1.2, are able to communicate unilaterally thanks to physical limitations in the construction of the circuit itself, that allows for current to only flow one way. This in turn leads to those pieces of the network being completely erasing any points of entry in the network.
There are similar software-side solutions to create unidirectional communication in protected networks, the most common of which is the use of firewalls, the main difference however between firewalls and data diodes is the level at which they operate. Taking into account the ISO/OSI stack, the conceptual standard for packet communications, firewalls most commonly operate around the network and transport layers, where TCP/IP operates. This means that attacks on the lower layers of the protocol (physical and data link) there are still vulnerabilities which could be attacked, layers which are protected by data diodes.
Given their security oriented nature, the main application of data diodes is in mission critical situations, where a part of the appliances has to be secured, but information about those appliances is still needed. Example use cases can be oil production plants, which are offshore, the status of which has to be monitored, and it is obviously not safe to connect directly to those plants as being on an open network could lead to attacks. Similarly air traffic and railway traffic monitoring appliances could be placed all over paths of interest.
Industrial grade data diodes however have a big drawback, which is their cost. A single diode device starts from 6000 dollars because of the complex circuits made for unidirectionality to work. Furthermore in industrial applications often a diode needs to be installed for each single piece of machinery, the cost of which can get very expensive. Luckily there are cheaper alternatives to industrial grade data diodes which do a similar job for most use case scenarios. For example a cheaper diode could be achieved using serial communication via an RS-232 cable. This works because the RS-232 standard for communication uses a pin to send data, and the other to receive data, unilateral communication thus could be achieved by cutting the receiving pin from the side that is sending the data, and similarly the transmitting pin could be cut from the side that receives the data. While this solution obviously does not function as smoothly as an industrially built diode, it will work for most applications.
Another issue with data diodes is that malfunctions in the private network could be difficult to signal to the end user. If the diode itself malfunctions for example, data will be lost and the machine itself has no way to diagnose the problem or to send a warning message. This problem could be solved with bidirectional communication, for example having the unprotected side ping the data diode periodically so as to check its status. However diodes are an unilateral communication device, so this procedure is not immediate. There is a way to make bidirectional communication possible, within a diode environment it is usually done by placing two data diodes, one from the protected network to the unprotected network (as seen in Figure 1.2), and another diode placed sending data the opposite way. In this bidirectional model data transfers can only be started from the protected side however, so as to take software precautions before starting communication [OWL18].




\chapter{Alstom group and their problem}

\section{Dissecting Alstom as a company}
\section{Alstom solution for railway track monitoring}
focusing on their shift to computer science oriented solutions

\chapter{Project}

\section{Alstom issues with interacting with machines under private networks}
The aim of my internship project with Alstom Group is to start by developing what in cybersecurity and IoT is called a data diode, a unidirectional communication device for data exchange.
The data that will be exchanged is diagnostic data caught by network monitoring program Zabbix. Zabbix will catch data from devices inside the protected network with a server running on a Raspberry Pi 3 and recreate the same Zabbix setup with the protected network values on the unprotected network. This Zabbix server will also run on a Raspberry Pi 3. These two devices are also the base behind the diode itself.
The importance behind this project is filling a market void created by industrial data diodes, which are very expensive devices. In use cases like this one, where the protected -> unprotected network links to make are plenty, implementing a high amount of them is a very costly move, but with software implementations (firewalls) or hardware implementations (serial or optical communication built to be unidirectional) it is possible to have cheaper data diodes that are almost as functional for most use case scenario. As mentioned earlier, the setup includes two Raspberry Pi model 3 computers, and a TTL-232R-3V3 cable with a specific configuration which we will highlight later.
\section{Goals of the thesis}
\section{Top-down description of the implemented solution}
As far as choosing the best Raspberry Pi model for the job, I have chosen to use the Raspberry Pi 3. For this type of project, finding a machine that can be used without cooling, that has enough RAM to handle possible big data chunks and that has a fast enough CPU at a low enough cost is key. Needing for it not to be cooled removes the Raspberry Pi 4 from the equation, since it prefers having a cooling solution and the increased clock speeds won't help with serial communication limitations, on the other hand the Raspberry Pi 2 already comes with 1Gb Ram, but misses the clock speed needed for multiple programs to run at a good enough speed.
This particular machine should be able to handle the data size we are planning to share (around 1Gb/day). The machines are configured as follows:
\begin{itemize}
    \item Raspberry Pi 3 Model B Ver1.2:
    \begin{itemize}
        \item Quad Core 1.2GHz Broadcom BCM2837 64bit CPU
        \item 1GB RAM
        \item BCM43438 wireless LAN and Bluetooth Low Energy (BLE) on board
        \item 100 Base Ethernet
        \item 40-pin extended pinout with GPIO support
        \item 4 USB 2 ports
        \item 4 Pole stereo output and composite video port
        \item Full size HDMI
        \item CSI camera port for connecting a Raspberry Pi camera
        \item DSI display port for connecting a Raspberry Pi touchscreen display
        \item Micro SD port for loading your operating system and storing data
        \item Upgraded switched Micro USB power source up to 2.5A
        \item Raspberry Pi OS Lite (latest version)
        \item 16GB microSD card
    \end{itemize}
\end{itemize}

Using the RPI-Imager tool I installed the latest Raspberry PI OS image. Upon first boot up, remember that RPi login is username: pi, password: raspberry.
As Raspberry Pi OS Lite is an OS without a desktop environment, we don't have a systems setting manager with a GUI, luckily we can use the utility \textit{raspi-config} to enable a series of features useful for testing software.
Using the command \textit{sudo raspi-config} we can open the utility, once in the main menu we change these specific settings:
\begin{itemize}
    \item In the system options:
    \begin{itemize}
        \item Enable WIFI for SSH purposes.
        \item Change username and password if needed.
    \end{itemize}
    \item In the interface options:
    \begin{itemize}
        \item Enable SSH. Using SSH is suggested for development on these units, as they are installs without a DE/WM(desktop environment/window manager) configuration, and thus it's best to use your preferred terminal emulator.
        \item Enable Serial interface, remembering to disable login shell but enable serial interface.
    \end{itemize}
    \item Disable bluetooth:
    \begin{itemize}
        \item We won't be using bluetooth for the project, so we disable it, making the final product more secure.
        \item We disable bluetooth by adding \textit{dtoverlay=disable-bt} in \textit{/boot/config.txt}.
    \end{itemize}
    \item Troubleshooting
    \begin{itemize}
        \item Check that UART is enabled in \textit{/boot/config.txt} by adding \textit{enable\_UART=1}.
        \item Check that serial console is disabled by removing \textit{console=serial0,115200} or \textit{console=ttyS0,115200} in \textit{/boot/cmdline.txt}.
        \item Check for permissions in the dialout group.
        \item Check that \textit{/dev/serial0} doesn't have a getty console running on it. In case it does, it can be disabled by using the commands: \textit{sudo systemctl stop serial-getty@ttyS0.service} and \textit{sudo systemctl disable serial-getty@ttyS0.service}.
    \end{itemize}
    \item Additional software:
    \begin{itemize}
        \item Run \textit{sudo apt-get install python3} to install Python in it's latest version.
        \item Run \textit{sudo apt-get install python3-pip} to install Python's package manager.
        \item Run {pip install pyserial} to install the serial library that we will be using to have the machines comunicate.
    \end{itemize}
    \item Run \textit{sudo apt-get upgrade} to make sure all packages are up to date.
\end{itemize}

\subsection{Serial cable}
As mentioned earlier, the serial communication will go through a TTL-232R-3V3 cable, which is USB-to-serial, with +3.3V TTL levels UART signals. The cable has 6-pins on one end, and USB on the other.
With the cable plugged in we run the command \textit{ls /dev/tty*} on both RPi machines so we have an idea of what port we use to send or receive data on each RPi in the serial communication, for the port that's sending data it will be \textit{/dev/ttyS0}, while for the other port it should be \textit{/dev/ttyUSB0}. Raspian will always map the serial port to the alias \textit{/dev/serial0}, so it is suggested to use the alias in programming as it doesn't differentiate between Raspberry Pi Models. Here below in Figure 1 there is an image showing how the cable is meant to be plugged in, while in the Appendix 1 there is the full extract from the documentation of the cable, showing the functionality of all the cable pins in detail.

\begin{figure}[h!]
    \centering
    \includegraphics[scale = 0.5]{img2}
    \caption{Pinout diagram}
    \label{fig:figure1}
\end{figure}

The cable pinout is peculiar as most pins are not used. The pins that get used are the ground pin and the pin that sends data over to the other RPi machine. Yellow pin is plugged in GPIO14, ground can be plugged in any of the RPi ground pins.

\subsection{How is it a diode?}
The importance behind this project is creating a diode-type communication, thus a unilateral comunication from a protected network to less protected or unprotected network. Below shown in Figure 2 how the two RPi interface within the network of Alstom devices, that is the first piece that shows how this setup has diode communication. The second piece is the pinout shown in Figure 1 in the chapter above, we can see that in the RPi placed in the private network, the serial pins that are plugged in are GND and TX, thus no data will be received through the cable by the RPi in the private network (RX pin is not plugged). These two factors create unidirectional communication similar to that of a diode.
\begin{figure}[h!]
    \includegraphics[scale = 1]{img3}
    \caption{Network diagram}
    \label{fig:figure2}
\end{figure}


\section{Program versions}
\subsection{Diagnostic test}
The objective of this test is to check if the two RPi are connected correctly. In Appendix 2 we have example code and example output that can be used to troubleshoot and check that the cable has been plugged in correctly.
In case this program won't run correctly, refer back to the troubleshooting section earlier. Make sure to run the code \textit{sndtest.py} on the machine where the GPIO pins are being used, and \textit{rcvtest.py} has to be run on the RPi where the USB is plugged in.
The scripts are started by running the command \textit{python3 rcvtest.py}, and then run \textit{python3 sndtest.py} on each machine.

\subsection{Data diode 0.1}
In this initial implementation we switch from sending text to sending an actual file via serial communication. The improvements from this version to the next will be improved baudrate, meaning faster file transfer.
In the next version SHA256 checksum will also be implemented to make sure file transfer is working correctly. Future considered improvements are to include a handshake system via CTS/RTS signals and cleaner code in general. The code can be found in Appendix 3.

\subsection{Data diode 0.2}
This version of the program has SHA256 integration to validate file transfer and improved transfer speeds thanks to a much better baudrate.
The limitations with this version of the program is that because of limitations of USB2.0 and limitations of the UART communication protocol, 115200 baudrate seems to be the maximum baudrate for communication without introducing errors, meaning that improving transfer speeds can't be done by increasing the baudrate, but improvements have to be found in other points of the code.
A comparative of baudrates can be seen in Table 1, where examples of how increasing baudrates helps with transfer speeds, and also a small example of whether calculating SHA256 while sending the file slows down communication, it will still be wise to split file transfer and calculations. Another important detail to note from the table is the approximate time for 1Gb file transfer, which given the current times might be a little too slow. In appendix 4 there are code examples for this version of the program as well.

\begin{table}[h!]
    \begin{tabular}{@{}|l|l|l|@{}}
        \toprule
        File size (mb) & Time (s) & SHA256 \\ \midrule
        1(9600baud)    & 1065.9   & Yes    \\ \midrule
        1              & 88.8     & No     \\ \midrule
        1              & 88.8     & Yes    \\ \midrule
        5              & 448      & Yes    \\ \midrule
        10             & 892.7    & Yes    \\ \midrule
        20             & 1785.4   & Yes    \\ \midrule
        1000           & ~89250   & Yes    \\ \bottomrule
    \end{tabular}
    \caption{Transfer times comparative}
    \label{tab:table 1}
\end{table}

\subsection{Data diode 0.3}
The aim with version 0.3 of the program is to initially format the code so as to make it much more readable. As far as technical improvements, this code had the massive flaw of wanting to send any type of file, but but using a transfer protocol built for text files, while at the same time opening the files in byte read as that's the only way of sending data through the serial port. Thus we will include byte read with a chunk size to be defined, intead of \textit{readline()}, so as to make less calls to the read and write methods, also we will be splitting up the SHA256 algorithm from the copying of file contents. Program will be upgraded to version 1.0 when it will be fully optimized for the final app. Code is in the appendix.


\section{Network monitoring}
\subsection{Why do we need network monitoring}
Having created a working data diode clone, we move on to the network monitoring part of the project. We need network monitoring software, so as to have a way to create the data that we then transfer via diode. Obviously the data will be, in the final version, a breakdown of the private network diagnostic values caught by the network monitor we choose. This data will be used to create a simulated version of the private network on the public network Raspberry Pi.
We will have thus to pick a network manager able to export and import data. \textit{PRTG Network Manager} is considered the best tool for this type of job, but it is only made to run on Windows server, and it is a closed source solution, so ARM based open source solution is preferred.
Well known market solutions for network monitoring are the following:
\begin{itemize}
    \item Cacti
    \item Nagios
    \item Zabbix
\end{itemize}
These solutions are all open source, all working on ARM and are all based on a LAMP stack. LAMP is an acronym for Linux, Apache, MySQL and PHP, but these programs are all flexible enough to be able to interchange Apache with Nginx or MySQL with database management systems such as PostgreSQL.

\subsubsection{Comparison between network monitoring software}
Starting from the premise that all these softwares allow for export of configuration and data, in picking the best network monitoring software for the job I will also have to evaluate other factors including the feasability of recreating the simulation of the protected network on the unprotected network,
Starting with Cacti, it's a stable solution since it has been in development since 2001. This leads sadly to it having less modern features compared to some of the other network managers, and this lack of features leads to Cacti being used less in the user space. Even if some of these features won't be used for our project, picking a program that has lots of users will make debugging for errors easier, as the community will have much more people around it. Cacti has a very straightforward way of importing data into it, but it has very little documentation on the idea of recreating a device which is not actually on the current network, so I decided to look at other options.
Nagios is the second options that I analyzed, it is stable, widely used, modern and has a big user space, sadly it has a very limited free version, and it is very comparable to Zabbix, which instead is free. Zabbix is the option I ended up going with for all the reasons mentioned above inlcuing some other factors such as the fact that their Raspberry Pi install guide is also very detailed, it includes the utility \textit{zabbix\_sender} which we can use to import data, and has another fan-made utility called \textit{Zabbix Agent Simulator} which we can use to create devices on a network where we don't have one. Being able to do exactly that, and import our data from the private network on those devices should give us the exact solution we are looking for.

\subsection{How does Zabbix work?}
Zabbix as a network monitoring tool has a series of main components, these are hosts (or agents), items, triggers.
Hosts are the devices which we want to monitor,

\subsection{Installation guide}
Initially I'll install Zabbix on the private Raspberry Pi for testing purposes, after having a working Zabbix on the private Raspberry Pi, I will work with the API and the other tools I mentioned to simulate setups. The install guide I followed is the one on the \textit{zabbix.com/download} websites, which has a detailed guide for exactly our RPi setup. Beware that the MySQL version in the install guide could be outdated, so I suggest updating to the newest version running the command \textit{sudo apt-get install mysql} before creating the database host. I also suggest \textit{mysql\_secure\_installation} as away to correctly set-up the MySQL installation on the machine.


\chapter{Conclusion}

\bibliographystyle{plain}

\chapter{Appendix}
\subsection{Appendix 1: Documentation extract}
\includegraphics[scale = 0.65]{img1}
\pagebreak
\subsection{Appendix 2: Diagnostic output and code example}
\lstinputlisting[language=Python, caption=sndtest.py]{./DiagnosticTest/sndtest.py}
\lstinputlisting[language=Python, caption=rcvtest.py]{./DiagnosticTest/rcvtest.py}
\lstinputlisting[caption=output.txt]{./DiagnosticTest/output.txt}
\pagebreak
\subsection{Appendix 3: Data diode 0.1 code example}
\lstinputlisting[language=Python, caption=sndfile0.1.py]{./1FileTransfer/sndfile.py}
\lstinputlisting[language=Python, caption=rcvfile0.1.py]{./1FileTransfer/rcvfile.py}
\pagebreak
\subsection{Appendix 4: Data diode 0.2 code example}
\lstinputlisting[language=Python, caption=sndfile0.2.py]{./2FileTransfer/sndfile.py}
\lstinputlisting[language=Python, caption=rcvfile0.2.py]{./2FileTransfer/rcvfile.py}

\chapter{Ringraziamenti}

\end{document} 
